\documentclass[9.5pt,a4paper]{article}

\usepackage[utf8]{inputenc}
\usepackage[T1]{fontenc}
\usepackage[french]{babel}
\usepackage[a4paper,left=1.15cm,right=1.15cm,top=0.65cm,bottom=0.45cm]{geometry}
\usepackage{enumitem}
\usepackage{titlesec}
\usepackage{xcolor}
\usepackage[hidelinks]{hyperref}

\pagestyle{empty}
\setlength{\parindent}{0pt}
\setlength{\parskip}{1.5pt}
\linespread{1.0}

\definecolor{Primary}{HTML}{1F2937}

\titleformat{\section}{\large\bfseries\color{Primary}}{}{0em}{}[\titlerule]
\titlespacing{\section}{0pt}{4pt}{2.5pt}

\setlist[itemize]{leftmargin=1.2em, itemsep=3pt, topsep=2pt, parsep=0pt}

\newcommand{\cvitem}[2]{\textbf{#1}: #2\\[2pt]}
\newcommand{\jobtitle}[3]{\textbf{#1}\hfill{\small #2}\\[1pt]\textit{#3}}

\begin{document}

\begin{center}
{\Huge\bfseries Ramy Lazghab}

\vspace{1.2mm}
{\large Ing\'enieur IA \& Machine Learning}

\vspace{1.2mm}
Paris, France $\bullet$ Ouvert \`a la mobilit\'e
\hspace{0.3cm}$\bullet$\hspace{0.3cm}
\href{mailto:ramy.lazghab@dauphine.eu}{ramy.lazghab@dauphine.eu}
\hspace{0.3cm}$\bullet$\hspace{0.3cm}
+33 7 59 68 79 72

\vspace{1.2mm}
\href{https://linkedin.com/in/ramy-lazghab-1464a8201}{LinkedIn}
\hspace{0.3cm}|\hspace{0.3cm}
\href{https://github.com/Rblaze23}{GitHub}
\hspace{0.3cm}|\hspace{0.3cm}
\href{https://rblaze23.github.io/RamyLazghab/}{Portfolio}

\end{center}

\section*{Profil}
Ing\'enieur IA concevant des syst\`emes LLM, agentiques, RAG et de ML pr\'edictif, du prototype \`a la mise en production, avec un fort accent sur l'\'evaluation et la fiabilit\'e. En fin de M2 Science des Donn\'ees \& IA \`a l'Universit\'e Paris Dauphine -- PSL.

\section*{Comp\'etences techniques}
\cvitem{Programmation}{Python, SQL, Java, JavaScript, C, R}
\cvitem{Machine Learning}{Deep learning, apprentissage par renforcement, NLP (traitement du langage naturel), transformers, fine-tuning, gradient boosting, explicabilit\'e des mod\`eles (SHAP) ; PyTorch, TensorFlow, XGBoost, LightGBM}
\cvitem{IA g\'en\'erative \& LLM}{RAG, syst\`emes multi-agents, \'evaluation de LLM, prompt engineering ; LangChain, LangGraph, Hugging Face, Anthropic API, OpenAI API, MCP, LangSmith}
\cvitem{MLOps, Cloud \& Donn\'ees}{Docker, Git, MLflow, GCP (Vertex AI, BigQuery, Cloud Run), API REST, PostgreSQL, Neo4j, bases vectorielles (FAISS, Qdrant), Spark, Pandas, NumPy, Streamlit}

\section*{Exp\'erience professionnelle}

\jobtitle{Ing\'enieur IA, stage de fin d'\'etudes}{Avr. 2026 -- aujourd'hui}{Relay X, Paris}

\begin{itemize}
\item \textbf{Oracle AI, pr\'evision du risque r\'eglementaire :} pr\'edit si un r\'egulateur restreindra un ingr\'edient sous 6 ou 18 mois. 1{\,}222 avis scientifiques structur\'es en un vecteur de 70 variables ; ensemble XGBoost + LightGBM, calibration isotonique, explications SHAP. \textbf{AUC-ROC 0,948 \`a 6 mois et 0,914 \`a 18 mois}, avec validation temporelle stricte sur avis non vus.
\item Enrichissement en temps r\'eel via un pipeline LangGraph (Anthropic Claude) \`a trois n\oe uds parall\`eles (PubMed, PubChem, web), restitu\'e en rapport Streamlit ; tra\c{c}age LangSmith.
\item \textbf{PIF (g\'en\'eration de dossiers) :} pipeline LLM \`a six agents produisant des Dossiers d'Information Produit conformes au R\`eglement (CE) 1223/2009, avec sections pilot\'ees par template et validation humaine de l'\'evaluation de s\'ecurit\'e. \textbf{Mis en production} sur la plateforme de production (Next.js, TypeScript, Express, Prisma, PostgreSQL).
\item Conception d'un benchmark de v\'erit\'e terrain (5 produits, 203 champs), boucle de correction sur trois versions : \textbf{affirmations non sourc\'ees r\'eduites de 87,5\,\% \`a 35,0\,\%, F1 port\'e de 81,8 \`a 89,5, pour 0,13\,\$ par dossier}.
\item \textbf{Magnet (radar r\'eglementaire) :} cascade de recherche \`a cinq niveaux (SQL d'abord, LLM si \'echec), rapports sourc\'es en trois passes. \'Evalu\'e sur 35 rapports g\'en\'er\'es \`a \textbf{88,6\,\% de tra\c{c}abilit\'e des sources} ; 75 tests de s\'ecurit\'e.
\end{itemize}

\vspace{1.8mm}
\jobtitle{D\'eveloppeur logiciel, stage}{F\'evr. -- Juin 2024}{Infinity Management}\\[2pt]
Plateforme e-learning personnalis\'ee par NLP (analyse de sentiment BERT) ; Django REST, React, Docker.

\section*{Formation}
\jobtitle{M2 Science des Donn\'ees \& Intelligence Artificielle}{2024 -- 2026 (en cours)}{Universit\'e Paris Dauphine -- PSL}

\vspace{1.8mm}
\jobtitle{Licence en G\'enie Informatique (IoT)}{2021 -- 2024}{Facult\'e des Sciences de Tunis}

\section*{Projets s\'electionn\'es}
\begin{itemize}
\item \textbf{TelecomPlus :} support client multi-agents (RAG + routage SQL), \'evaluation LLM-as-a-Judge, tra\c{c}age Langfuse \& MLflow. \textit{LangChain, LangGraph, FAISS, PostgreSQL}
\item \textbf{RAGenius :} assistant IA de questions-r\'eponses s\'emantiques sur PDF et CSV. \textit{LangChain, Hugging Face, FAISS}
\item \textbf{Movie Recommender :} syst\`eme de recommandation de bout en bout sur GCP, expos\'e en API REST. \textit{BigQuery ML, Vertex AI, Cloud Run}
\item \textbf{Pr\'ediction d'Alzheimer :} ensemble XGBoost et deep learning pour la d\'etection pr\'ecoce ; SHAP pour l'interpr\'etabilit\'e clinique. \textit{PyTorch, XGBoost, SHAP}
\end{itemize}

\section*{Certifications \& Distinctions}
\begin{itemize}
\item \textbf{Certifications :} \href{https://academy.langchain.com/certificates/nqrsewnhol}{LangChain Academy} ; \href{https://www.credly.com/badges/9ec1dc8b-494f-42ee-a314-ba30b40342de/public_url}{Google Cloud, BigQuery ML}
\item \textbf{3\textsuperscript{e} place}, RAISE Summit AI Hackathon Paris ; \textbf{Laur\'eat}, Hack for Good (MoodSync) ; EY Hack for Smart Insurance ; IEEE Xtreme 15.0/16.0 ; membre IEEE, ITLAB \& AIESEC
\item \textbf{Langues :} Anglais, courant (IELTS) $\vert$ Fran\c{c}ais, professionnel
\end{itemize}

\end{document}
